%%%%%%%%%%%%%%%%%%%%%%%%%%%%% BLOCK: ps_results %%%%%%%%%%%%%%%%%%%%%%%%%%%%%%%
\subsection*{Results}

Table~\ref{tab:subsources} reports $\Delta_{AB}$ for the fifteen units for
which both prior modes were run.  Four are data-driven: D~source
($\Delta_{AB} = -9$~yr), D~Frame ($-20$~yr), Song of the Sea ($+21$~yr),
and the Jeremiah oracle holdout ($-27$~yr).  Eight are mildly
prior-influenced, including the P source ($+39$~yr).  Three are
prior-dominated: the JE source ($+139$~yr), Habakkuk ($-138$~yr) and
Daniel ($-98$~yr).

Two features of this pattern deserve comment.  First, the units classified
as prior-dominated include two of the four held-out validation texts.  A
validation design in which the held-out units are the ones whose answers
are most determined by the prior is not a test of the model, and this
observation was the entry point to the leakage analysis reported under
\nameref{sec:validation}.  Second, the classification is not a function of
text length: Daniel (3{,}437 words) and Habakkuk (897 words) are
prior-dominated, while the Song of the Sea (433 words) is data-driven.  It
is a function of prior width.  Units carrying a tight scholarly $\sigma$
move little when the prior is exchanged, not because the evidence is
strong but because the evidence is being overridden.

We therefore report $\Delta_{AB}$ as a diagnostic of prior influence rather
than as a robustness certificate.  A small $\Delta_{AB}$ means only that
the two priors agree; combined with the precision decomposition of
Table~\ref{tab:leakage}, it does not by itself establish that a date is
driven by linguistic evidence.

%%%%%%%%%%%%%%%%%%%%%%%%%%%% BLOCK: lbh_score %%%%%%%%%%%%%%%%%%%%%%%%%%%%%%%%%
\subsection*{Model-independent feature scoring}

It is useful to have a description of each unit that can be checked against
the \bhsa{} database without running any model.  Each feature is normalized
to a $[0,1]$ scale by reference to the fitted trend at two anchor dates:
%
\begin{equation}
  s_j = \frac{x_j - \hat{f}_j(d_{\text{CBH}})}{
    \hat{f}_j(d_{\text{LBH}}) - \hat{f}_j(d_{\text{CBH}})},
  \label{eq:lbh_score}
\end{equation}
%
where $d_{\text{CBH}} = 720$~BCE and $d_{\text{LBH}} = 250$~BCE, so that
$s_j = 0$ is \cbh{}-like, $s_j = 1$ is \lbh{}-like, and values outside
$[0,1]$ lie beyond the anchors in either direction.  The mean score
$\bar{s}$ summarises a unit.

On the present corpus the source composites order
D ($\bar{s} = -0.30$) $<$ P ($+0.15$) $<$ JE ($+0.53$), with the
Deuteronomic legal code the most \cbh{}-like unit in the study
($\bar{s} = -0.63$) and the Song of the Sea close to the \cbh{} anchor
($\bar{s} = -0.00$).  We report these figures because S3~Table lists the
individual features driving them, and because a reader who distrusts the
model can inspect them directly.  We caution against reading the ordering
as a chronology.  $\bar{s}$ is an unweighted average over features selected
for their correlation with date, and it makes no allowance for the genre
composition of a unit; D's legal core is formulaic and archaizing in
diction, which is sufficient to explain its position without any claim
about its date.  Earlier versions of this work reported a different
ordering (D $<$ JE $<$ P) from an earlier corpus and offered a
substantive interpretation of it.  That interpretation does not survive
the corpus revision, and the instability of the ordering across corpus
versions is itself a reason to treat $\bar{s}$ as descriptive rather than
probative.
